\pg{437} a existencia de uma Intelligencia soberana e infinita. Na verdade, todo o agente \emph{irracional}, nas suas operações, tende para um fim, empregando para isso os meios mais efficazes, e executando um plano com maravilhosa exactidão, constancia e uniformidade. Este fim, esta escolha de meios e este plano, devem ser effeitos de uma intelligencia \emph{infinita}; porque só uma intelligencia \emph{infinita} podia determinar a cada agente \emph{irracional} um fim proporcionado, modelar planos complicados e ao mesmo tempo perfeitos, e, para os realizar, servir-se da forças cegas da natureza. A intelligencia \emph{infinita} é propria de \textsc{Deus}. Logo existe \textsc{Deus} \footnote{Tambem este ponto ácerca da \emph{finalidade} nas operações dos agentes \emph{irracionaes} foi sufficientemente tratado na \emph{Cosmologia geral}.}.

14. \textbf{Argumentos moraes.} --- Os argumentos \emph{moraes} fundam-se em que, --- \emph{a}) se existe uma lei \emph{moral}, deve existir um legislador, --- \emph{b}) se todos os povos admittiram a existencia de um Ente supremo, Este deve existir.

a) A consciencia attesta-nos que temos gravada no nosso coração uma lei moral, á qual não nos podemos subtrahir, e que o cumprimento d'essa lei produz tranquillidade, assim como a sua transgressão causa remorso. Ora a existencia d'esta lei, como a respectiva sancção, suppõe a existencia de um legislador. Este legislador não póde ser o proprio homem; aliás seria inexplicavel o sentimento da nossa dependencia de um poder superior. Logo deve ser um Ente superior ao homem e dotado de virtude ou poder \emph{infinito}; porque só um ente infinitamente poderoso podia gravar no coração dos homens essa lei. O Ente superior, de que o homem depende e a que deve obedecer, é \textsc{Deus}. Logo existe \textsc{Deus} \footnote{Este argumento, apresentado tambem por \emph{Kant} para a demonstração da existencia de \textsc{Deus}, será mais desenvolvido na \emph{Moral}.\par Aqui apenas advertimos que a existencia da \emph{lei moral} considera-se como um \emph{facto experimental}, attestado pela consciencia. Provado o facto da existencia da \emph{lei}, legitimamente deduz-se a existencia de um \emph{legislador}.}.
